%==============================================================
%  Yifan Yang — Academic Curriculum Vitae
%  Overleaf-ready.  Compile with pdfLaTeX (or XeLaTeX).
%  Style: classic academic (serif, ruled section headings)
%==============================================================
\documentclass[10pt,letterpaper]{article}

%---------------- Page geometry -------------------------------
\usepackage[left=0.75in,right=0.75in,top=0.70in,bottom=0.70in,footskip=0.3in]{geometry}

%---------------- Encoding & fonts ----------------------------
\usepackage[utf8]{inputenc}
\usepackage[T1]{fontenc}
\usepackage{charter}          % Bitstream Charter: clean academic serif
\usepackage{textcomp}
\usepackage{microtype}

%---------------- Layout packages -----------------------------
\usepackage{titlesec}
\usepackage{enumitem}
\usepackage{xcolor}
\usepackage{ragged2e}
\usepackage[hidelinks]{hyperref}
\usepackage{fancyhdr}

%---------------- Colours -------------------------------------
\definecolor{accent}{HTML}{500000}   % Texas A&M maroon
\definecolor{rulegray}{HTML}{999999}

\hypersetup{
  colorlinks=true,
  linkcolor=accent,
  urlcolor=accent,
  citecolor=accent,
  pdftitle={Yifan Yang -- Curriculum Vitae},
  pdfauthor={Yifan Yang}
}

%---------------- Section formatting --------------------------
\titleformat{\section}
  {\normalfont\large\bfseries\color{accent}\scshape}
  {}{0em}{}[\vspace{-0.15em}{\color{rulegray}\titlerule[0.7pt]}]
\titlespacing*{\section}{0pt}{1.15em}{0.55em}

\titleformat{\subsection}
  {\normalfont\normalsize\bfseries}
  {}{0em}{}
\titlespacing*{\subsection}{0pt}{0.75em}{0.25em}

\titleformat{\subsubsection}[runin]
  {\normalfont\normalsize\itshape}
  {}{0em}{}[.]
\titlespacing*{\subsubsection}{0pt}{0.5em}{0.4em}

%---------------- Footer --------------------------------------
\pagestyle{fancy}
\fancyhf{}
\renewcommand{\headrulewidth}{0pt}
\fancyfoot[L]{\footnotesize\itshape Yifan Yang}
\fancyfoot[R]{\footnotesize\thepage}

%---------------- Helper macros -------------------------------
% \entry{Title / Institution}{Right-hand info (dates)}
\newcommand{\entry}[2]{%
  \par\vspace{0.30em}\noindent\textbf{#1}\hfill{\small #2}\par}

% \subentry{Role / degree}{Right-hand info (location)}
\newcommand{\subentry}[2]{%
  \par\nobreak\vspace{0.05em}\noindent #1\hfill{\small\itshape #2}\par}

% Compact bullet list
\newenvironment{tight}
  {\begin{itemize}[leftmargin=1.25em,itemsep=0.10em,topsep=0.20em,parsep=0pt,label=\textbullet]}
  {\end{itemize}}

% Compact list with no bullet marker (for award / service lines)
\newenvironment{plainlist}
  {\begin{itemize}[leftmargin=1.1em,itemsep=0.10em,topsep=0.20em,parsep=0pt,label={}]}
  {\end{itemize}}

% Numbered-reference list where numbers are typed manually
\newenvironment{reflist}
  {\begin{list}{}{%
     \setlength{\leftmargin}{1.7em}\setlength{\itemindent}{-1.7em}%
     \setlength{\itemsep}{0.22em}\setlength{\parsep}{0pt}\setlength{\topsep}{0.25em}}}
  {\end{list}}

\newcommand{\award}[1]{\textbf{\textit{#1}}}
\newcommand{\dt}[1]{\textbf{#1}\hspace{0.6em}}

\setlength{\parindent}{0pt}
\raggedright

%==============================================================
\begin{document}
%==============================================================

%---------------- Header --------------------------------------
\begin{center}
  {\LARGE\bfseries\scshape Yifan Yang}\\[0.45em]
  {\small
   College Station, TX \,$\cdot$\,
   \href{mailto:yyf990925@gmail.com}{yyf990925@gmail.com} \,$\cdot$\,
   (213)\,590-8972 \,$\cdot$\,
   \href{https://rayford295.github.io/}{rayford295.github.io} \,$\cdot$\,
   \href{https://scholar.google.com/citations?user=B-fiSHwAAAAJ}{Google Scholar}}
\end{center}
\vspace{0.35em}

%==============================================================
\section{Education}
%==============================================================

\entry{Texas A\&M University}{2024 -- present}
\subentry{Ph.D. Student, Geography (GIScience) \, --- \, GPA 3.80/4}{College Station, TX, USA}
\subentry{Advisor: Dr.~Lei Zou}{}

\entry{University of Southern California}{2022 -- 2024}
\subentry{Master of Science, Spatial Data Science \, --- \, GPA 3.96/4}{Los Angeles, CA, USA}
\subentry{Advisors: Dr.~John P. Wilson and Dr.~Siqin Wang}{}

\entry{Hainan University}{2018 -- 2022}
\subentry{Bachelor of Engineering, Software Engineering \, --- \, Average score 87.4/100}{Haikou, China}
\subentry{Advisors: Dr.~Jieren Cheng and Dr.~Qi Qi}{}

%==============================================================
\section{Research Interests}
%==============================================================

Geographic Information Science (GIScience) \,$\cdot$\, Geospatial Artificial Intelligence (GeoAI) \,$\cdot$\, Cartography \,$\cdot$\, Spatial Data Science \,$\cdot$\, Autonomous GeoAI \,$\cdot$\, Generative AI \,$\cdot$\, Multimodal AI \,$\cdot$\, Disaster Resilience

\vspace{0.5em}
\textit{Vision:} To transcend the boundaries of screens and make the real world our playground of intelligence --- where AI, space, and humanity coexist and co-create.

%==============================================================
\section{Current Projects}
%==============================================================

\entry{Sat2Street-DisasterGen --- Satellite-to-Street View Synthesis for Post-Disaster Assessment}{\href{https://github.com/rayford295/Sat2Street-DisasterGen}{GitHub}}
\vspace{0.15em}
Developed a multimodal generative framework for satellite-to-street view synthesis under post-disaster scenarios. The system benchmarks cross-view image generation strategies that integrate diffusion models, ControlNet, and VLM-guided semantic conditioning to synthesize realistic street-view imagery from post-disaster satellite imagery.
\begin{tight}
  \item Designed a unified evaluation framework ensuring consistent satellite inputs, data splits, and protocols across generative methods
  \item Implemented multiple cross-view generation pipelines, including Pix2Pix (cGAN), SD1.5 + ControlNet, VLM-guided synthesis, and Disaster-MoE (Mixture-of-Experts)
  \item Built a multi-stage LoRA training pipeline for damage-severity-aware generation (Mild / Moderate / Severe)
  \item Evaluated outputs using pixel-level metrics (FID, SSIM, LPIPS), semantic consistency (ResNet-based classification), and VLM perceptual scoring
  \item Developed reproducible dataset structure, inference scripts, and evaluation modules for controlled benchmarking
\end{tight}

\entry{RAPID --- A Reproducible Multi-Agent Pipeline for Interpretable Disaster Damage Assessment from Satellite and Street-View Imagery}{\href{https://github.com/rayford295/RAPID}{GitHub}}
\vspace{0.15em}
Developed an autonomous multi-agent GeoAI framework for hyperlocal, interpretable, and near--real-time disaster damage assessment. The system orchestrates multiple specialized agents leveraging vision--language foundation models to perform cross-view perception, damage recognition, reasoning, and automated report generation without task-specific retraining.
\begin{tight}
  \item Designed a modular multi-agent architecture, including Disaster Perception, Image Restoration, Damage Recognition, and Disaster Reasoning agents
  \item Integrated multimodal inputs: satellite imagery, street-view imagery, textual cues, and temporal change signals
  \item Implemented agent-based orchestration for zero-/few-shot disaster analysis using VLMs (e.g., GPT, Gemini)
  \item Enabled cross-view semantic alignment and causal reasoning for damage assessment and recovery recommendation
  \item Built an interpretable evaluation and reporting pipeline for structured disaster intelligence outputs
\end{tight}

\entry{DamageArbiter --- CLIP-Enhanced Multimodal Framework for Hurricane Damage Assessment}{\href{https://github.com/rayford295/DamageArbiter}{GitHub}}
\vspace{0.15em}
Developed a CLIP-enhanced multimodal arbitration framework for hyperlocal hurricane damage assessment using bi-temporal street-view imagery. The system integrates vision-language models (e.g., CLIP) with deep vision backbones (e.g., ViT-B/16) to improve the robustness of damage classification under real-world post-disaster uncertainty.
\begin{tight}
  \item Designed a multimodal arbitration mechanism to reconcile conflicting visual and semantic signals
  \item Implemented image classification pipelines using ViT-B/16 and transformer-based architectures
  \item Evaluated performance on hurricane street-view datasets with precision, recall, F1, and cross-validation
  \item Built a reproducible codebase, including dataset structure, model training scripts, and figure generation
\end{tight}

\entry{BiTemporal-StreetView --- Hyperlocal Disaster Assessment via Bi-Temporal Street-View Imagery}{\href{https://github.com/rayford295/Bi-Temporal-StreetView}{GitHub}}
\vspace{0.15em}
Developed a bi-temporal street-view image analysis framework for hyperlocal disaster damage assessment using pre- and post-disaster imagery. The system employs dual-channel architectures with pre-trained vision backbones (Swin Transformer, ConvNeXt) to improve classification accuracy and the interpretability of damage-severity estimation.
\begin{tight}
  \item Designed a dual-channel fusion architecture for comparative reasoning between pre- and post-disaster imagery
  \item Constructed a dataset of 2,249 labeled street-view image pairs with fine-grained severity annotations
  \item Achieved performance improvement from 66.14\% (post-only) to 77.11\% (bi-temporal) classification accuracy
  \item Applied Grad-CAM visualization to demonstrate enhanced spatial attention using pre-disaster inputs
  \item Evaluated multiple fusion strategies (cross-attention, weighted fusion, Siamese difference)
  \item Enabled fine-grained damage mapping for climate-resilient urban planning
\end{tight}
\vspace{0.25em}
\textit{Industry engagement:} The research is promoted by \textbf{Mosaic --- Geospatial Imaging Leaders} (Prague, Czech Republic), highlighting its relevance to 360\textdegree{} street-level surveying, urban mapping, and disaster-intelligence applications. \href{https://www.mosaic51.com/featured/how-texas-am-is-pioneering-damage-assessment-using-mosaic-imagery-from-hurricane-milton/}{[feature article]}

%==============================================================
\section{Professional Experience}
%==============================================================

\entry{Graduate Research Assistant, Texas A\&M University}{Aug 2024 -- Present, TX}
\begin{tight}
  \item Conduct research on GeoAI, multimodal disaster assessment, and spatial intelligence using remote sensing and street-view imagery.
  \item Develop deep learning pipelines (CLIP, BLIP, ViT, ConvNeXt, Swin, LLM-based reasoning) for bi-temporal disaster perception and geospatial representation.
  \item Build geospatial databases, run spatial analysis (PostGIS, ArcGIS Pro), and implement Python ML workflows for classification, localization, and explainability (Grad-CAM, semantic probing).
\end{tight}

\entry{Research Intern, Spatial Data Lab, Harvard University}{Mar 2025 -- Aug 2025, Remote}
\begin{tight}
  \item Conducted research on wildfire modeling and Autonomous GeoAI, focusing on Los Angeles wildfire dynamics using multimodal geospatial data (satellite imagery, fire perimeters, VIIRS/GOES data, and weather layers).
  \item Developed RL- and POMDP-based frameworks to simulate wildfire spread and evaluate autonomous decision-making strategies for response, resource allocation, and risk-aware planning.
\end{tight}

\entry{Graduate Teaching Assistant, Texas A\&M University}{Aug 2024 -- Dec 2024, TX}
\begin{tight}
  \item Assisted in teaching Cartography (GEOG/GEOS course), including weekly labs, visualization exercises, and map design principles.
  \item Guided students in applying GIScience concepts, cartographic design, spatial thinking, and geovisualization techniques using ArcGIS Pro and related tools.
\end{tight}

\entry{Student Worker (Graduate Research), University of Southern California}{Sept 2023 -- May 2024, CA}
\begin{tight}
  \item Conducted research on urban tree shade modeling to evaluate microclimate impacts and heat mitigation benefits across Los Angeles neighborhoods.
  \item Built 3D urban environment models in ArcGIS Pro using building footprints, tree canopy data, LiDAR, and DSM/DEM layers.
  \item Performed sun--shade and solar radiation analysis to quantify shading patterns, assess thermal comfort variation, and support urban greening strategies.
\end{tight}

\entry{Data Engineering and AI Model Fine-tuning, Vitally AI}{Feb 2023 -- May 2023, Remote}
\begin{tight}
  \item Fine-tuned Stable Diffusion models using DreamBooth for personalized pet-style image generation, including embedding extraction, class/instance image balancing, and training with LoRA and prior-preservation loss to capture unique style features.
  \item Prepared and curated training datasets, performed image preprocessing (cropping, normalization, augmentations), optimized model configurations (learning rate, steps, scheduler, batch size).
  \item Developed Python scripts and automated data pipelines for dataset preparation, model training, and inference deployment; integrated workflows using PyTorch, HuggingFace Diffusers, and custom training utilities to support AI product prototyping and internal model experimentation.
\end{tight}

%==============================================================
\section{Publications}
%==============================================================

\href{https://scholar.google.com/citations?user=B-fiSHwAAAAJ}{Google Scholar Page} \,---\, Total citations: \textbf{382} \,$\cdot$\, \textit{h}-index: \textbf{6} \,$\cdot$\, \textit{i10}-index: \textbf{5}

\subsection{Publications as the First Author}

\subsubsection{Journal Articles}
\begin{reflist}
  \item[] [1] Yang, \textbf{Y.}, Zou, L., Zhou, B., Li, D., Lin, B., Abedin, J., \& Yang, M. (2025). Hyperlocal disaster damage assessment using bi-temporal street-view imagery and pre-trained vision models. \textit{Computers, Environment and Urban Systems}, 121, 102335.
  \item[] [2] Yang, \textbf{Y.}, Wang, S., Li, D., Sun, S., \& Wu, Q. (2024). GeoLocator: A Location-Integrated Large Multimodal Model (LMM) for Inferring Geo-Privacy. \textit{Applied Sciences}, 14(16), 7091.
\end{reflist}

\subsubsection{Book Chapters}
\begin{reflist}
  \item[] [1] Yang, \textbf{Y.}, \& Borrelli, D. (2025). Object detection and segmentation of trees using Text SAM in ArcGIS Online. In D. M. Ruddell \& D. Ter-Ghazaryan (Eds.), \textit{Security First: Geospatial Workflows for a Safe and Equitable World} (pp.~113--122). Redlands, CA: Esri Press.
\end{reflist}

\subsubsection{Conference Proceedings}
\begin{reflist}
  \item[] [1] Yang, \textbf{Y.}, \& Zou, L. (2025). Perceiving Multidimensional Disaster Damages from Street-View Images Using Visual-Language Models. \textit{Abstracts of the ICA}, 10, 310. \award{Best Student Paper Award}
  \item[] [2] Yang, \textbf{Y.}, Zou, L., \& Jepson, W. (2026). Satellite-to-Street: Synthesizing Post-Disaster Views from Satellite Imagery via Generative Vision Models. \textit{ArXiv}. \href{https://arxiv.org/abs/2603.20697}{arXiv:2603.20697}. \textit{Accepted by the 2026 IEEE International Geoscience and Remote Sensing Symposium.}
\end{reflist}

\subsubsection{Preprints}
\begin{reflist}
  \item[] [1] Yang, \textbf{Y.}, Gong, W., Zhang, K., Zou, L., Tu, Z., Li, H., Li, Z., \& Ye, X. (2026). RAPID: A Reproducible Multi-Agent Pipeline for Interpretable Disaster Damage Assessment from Satellite and Street-View Imagery. \textit{ArXiv}. \href{https://arxiv.org/abs/2606.21819}{arXiv:2606.21819}
  \item[] [2] Yang, \textbf{Y.}, Zou, L., Gong, W., Fu, K., Li, Z., Wang, S., \ldots\ \& Tian, H. (2026). DamageArbiter: A CLIP-Enhanced Multimodal Arbitration Framework for Hurricane Damage Assessment from Street-View Imagery. \textit{arXiv preprint} arXiv:2603.14837.
\end{reflist}

\subsection{Publications as a Co-Author}

\subsubsection{Journal Articles}
\begin{reflist}
  \item[] [1] Wang, S., Liu, H., Wu, C. Y., Huang, X., Wang, R., \textbf{Yang, Y.}, Corcoran, J., Lai, S., Xia, X., \& Liu, Y. (2026). The Boiling Frog Effect: Global Warming Delays Emotional Impacts of Air Pollution in Warmer Climates. \textit{Journal of Hazardous Materials}, 142440. \href{https://doi.org/10.1016/j.jhazmat.2026.142440}{doi:10.1016/j.jhazmat.2026.142440}
  \item[] [2] Li, Z., Li, H., \textbf{Yang, Y.}, Wang, S., \& Zhu, Y. (2026). Integrating earth observation data into the tri-environmental evaluation of the economic cost of natural disasters: A case study of 2025 LA wildfire. \textit{International Journal of Applied Earth Observation and Geoinformation}, 150, 105342. \href{https://doi.org/10.1016/j.jag.2026.105342}{doi:10.1016/j.jag.2026.105342}
  \item[] [3] Ye, X., Gong, W., \textbf{Yang, Y.}, Zou, L., Tu, Z., Huang, X., \ldots\ \& Wu, L. (2026). Towards Agentic Urban Digital Twins (AUDiTs): advancing new urban science through Human-AI co-learning agents. \textit{Urban Informatics}, 5(1), 9.
  \item[] [4] Wang, Y., Yin, H., \textbf{Yang, Y.}, Zhao, C., \& Wang, S. (2025). Navigating spatial inequities in freight truck crash severity via counterfactual inference in Los Angeles. \textit{Journal of Transport Geography}, 128, 104387.
  \item[] [5] Gong, W., Ye, X., Wu, K., Jamonnak, S., Zhang, W., \textbf{Yang, Y.}, \& Huang, X. (2025). Integrating spatiotemporal vision transformer into digital twins for high-resolution heat stress forecasting in campus environments. \textit{Journal of Planning Education and Research}, 1--15. \href{https://doi.org/10.1177/0739456X251391121}{doi:10.1177/0739456X251391121}
  \item[] [6] Zhang, Y., Chen, J., Liu, B., \textbf{Yang, Y.}, Li, H., Zheng, X., Chen, X., Ren, T., \& Xiong, N. (2020). COVID-19 public opinion and emotion monitoring system based on time series thermal new word mining. \textit{Computers, Materials \& Continua}, 64(3), 1415--1434. \href{https://doi.org/10.32604/cmc.2020.011316}{doi:10.32604/cmc.2020.011316}
  \item[] [7] Shi, C., Mou, G., Wang, H., Liu, T., \textbf{Yang, Y.}, \& Li, Z. (2020). Different Loss Functions Used in the Low-rank Approximation. \textit{International Core Journal of Engineering}, 6(11), 360--368.
\end{reflist}

\subsubsection{Conference Proceedings}
\begin{reflist}
  \item[] [1] Zou, L., Mandal, D., Zhou, B., \textbf{Yang, Y.}, \& Yang, M. (2025). Resilience in 4D: AI-Driven Geospatial Digital Twins for Urban Flood Simulation and Management. \textit{AGU25}.
  \item[] [2] Yan, B., Wang, J., Cheng, J., Zhou, Y., Zhang, Y., \textbf{Yang, Y.}, \ldots\ \& Liu, B. (2021, June). Experiments of federated learning for COVID-19 chest X-ray images. In \textit{International Conference on Artificial Intelligence and Security} (pp.~41--53). Cham: Springer International Publishing.
\end{reflist}

\subsubsection{Preprints}
\begin{reflist}
  \item[] [1] Gong, W., Srivastava, U., Wang, Y., Jia, Y., Wu, Q., Bai, W., \ldots\ \& Ye, X. (2026). Earth Embeddings Reveal Diverse Urban Signals from Space. \textit{arXiv preprint} arXiv:2604.03456.
  \item[] [2] Lin, B., Zou, L., Tian, H., Cai, H., \textbf{Yang, Y.}, \& Zhou, B. (2026). Predicting Healthcare System Visitation Flow by Integrating Hospital Attributes and Population Socioeconomics with Human Mobility Data. \textit{arXiv preprint} \href{https://arxiv.org/abs/2601.15977}{arXiv:2601.15977}.
  \item[] [3] Wu, Q., Xu, Y., Xiao, T., Xiao, Y., Li, Y., Wang, T., \ldots\ \& \textbf{Yang, Y.} (2024). Surveying attitudinal alignment between large language models vs. humans towards 17 sustainable development goals. \textit{arXiv preprint} arXiv:2404.13885.
\end{reflist}

%==============================================================
\section{Presentations (Presenting Author)}
%==============================================================

\begin{reflist}
  \item[] [1] \dt{2026} RAPIDMap: Rapid multi-Agent Pipeline for Interpretable Disaster Mapping from Satellite and Street-view Imagery. \textit{Cartography and Geographic Information Society (CaGIS) Conference}, St.~Louis, Missouri, September 8--11, 2026.
  \item[] [2] \dt{2026} DamageArbiter: A Disagreement-driven Arbitration Framework for Hurricane Damage Assessment from Street-View Imagery. \textit{AAG GISS Specialty Group Honors Competition for Student Papers 2}, AAG 2026 Annual Meeting, San Francisco, California.
  \item[] [3] \dt{2026} Satellite-to-Street: Synthesizing Post-Disaster Views from Satellite Imagery via Generative Vision Models. \textit{IEEE International Geoscience and Remote Sensing Symposium (IGARSS)}, Washington, D.C., August 13, 2026.
  \item[] [4] \dt{2026} Satellite-to-Street: Synthesizing Post-Disaster Views from Satellite Imagery via Generative Vision Models. \textit{Symposium on Spatiotemporal Data Science}, Virginia Tech Academic Building One, Alexandria, VA, August 7--8, 2026.
  \item[] [5] \dt{2025} Seeing Disagreements: An Explainable Multimodal Framework for Disaster Assessment and Spatial Mapping. \textit{Texas A\&M University GIS Day Student Paper Competition}, Texas A\&M University, College Station, USA (Oral).
  \item[] [6] \dt{2025} Yang, Y., Gong, W., Zhang, K. DisasterVLP: Perceiving Multidimensional Disaster Damages from Street-View Images Using Visual-Language Models. \textit{Summer School on Cyberinfrastructure and Disaster Resilience}, August 15, 2025, Texas A\&M University, College Station, USA (Oral).
  \item[] [7] \dt{2025} Perceiving Multidimensional Disaster Damages from Street-View Images Using Visual-Language Models. \textit{Student Paper Session, International Cartographic Conference (ICC)}, Vancouver, Canada.
  \item[] [8] \dt{2025} DisasterVLP: A Vision-Language Pretrained Framework for Multidimensional Disaster Damage Assessment Using Street-View Images. \textit{GeoAnalytics for Sustainable and Livable Cities Symposium, International Cartographic Conference (ICC)}, Vancouver, Canada.
  \item[] [9] \dt{2025} Hyperlocal Disaster Damage Assessment Using Bi-Temporal Street-View Imagery and Pre-Trained Image Processing Models. \textit{GISS Specialty Group Paper Competition II}, AAG 2025 Annual Meeting, Detroit, Michigan.
  \item[] [10] \dt{2024} GeoLocator: A Location-Integrated Large Multimodal Model for Inferring Geo-Privacy. \textit{Symposium on GeoAI and Deep Learning for Geospatial Research: Human-Centered Geospatial Data Science III}, AAG 2024 Annual Meeting, Honolulu, Hawai\textquoteleft i.
  \item[] [11] \dt{2024} GeoLocator: A Location-Integrated Large Multimodal Model (LMM) for Inferring Geo-Privacy. \textit{Spatial Data Science Symposium --- Thematic Session ``Geoprivacy Challenges and Solutions in the Digital Society.''} Online Symposium.
  \item[] [12] \dt{2024} GeoLocator: A Novel GeoAI Tool Making World Travel without Barriers \& Urban Tree Shade Model. \textit{2024 Los Angeles Geospatial Summit}, February 23, 2024, Los Angeles, California.
  \item[] [13] \dt{2024} GeoLocator: A Location-Integrated Large Multimodal Model for Inferring Geo-Privacy. \textit{AGI Leap Summit --- Multimodality Session}, SuperAGI, February 29, 2024 (Virtual).
\end{reflist}

%==============================================================
\section{Session Organizer \& Chair}
%==============================================================

\begin{reflist}
  \item[] [1] Spatiotemporal AI for Urban and Environmental Intelligence: Sensing, Inference, and Human-Centered Applications. \textit{Symposium on Spatiotemporal Data Science}, Virginia Tech Academic Building One, Alexandria, VA, August 7--8, 2026. \textit{Session Organizer.}
  \item[] [2] GeoAI and Data Science for Disaster Resilience. \textit{33rd International Conference on Geoinformatics (GeoInformatics 2026)}, Singapore. Co-organized with Lei Zou, Yi Qiang, Yingjie Hu, Qunying Huang, Xiao Huang, and Bing Zhou. \href{https://blog.nus.edu.sg/cpgis2026/programme/sessions/}{[programme]}
  \item[] [3] GeoAI and Deep Learning Symposium: GeoAI for Disaster Resilience. \textit{AAG 2026 Annual Meeting}, San Francisco, California. Co-organized with Lei Zou, Yingjie Hu, Qunying Huang, Yi Qiang, Marcela Suarez, and Morteza Karimzadeh.
  \item[] [4] Panel and Symposium Organization, \textit{GISER Symposium: Past, Present, and Future of GIScience}, AAG Annual Meeting, San Francisco, USA, 2026. \textit{Organizer.}
  \item[] [5] AAG Remote Sensing Specialty Group Student Illustrated Paper Competition. \textit{AAG 2026 Annual Meeting}, San Francisco, California. Co-organized with Ying Lu.
  \item[] [6] Symposium on Human Dynamics Research: Urban Environmental Intelligence and Human--Climate Interactions (Sessions 1--2). \textit{AAG 2025 Annual Meeting}, Detroit, Michigan. Co-organized with Wenjing Gong, Sisi Wang, Xinyue Ye, and Xiao Huang.
\end{reflist}

%==============================================================
\section{Selected Honors, Awards, and Sponsorships}
%==============================================================

\begin{plainlist}
  \item \dt{2026} Cartography and Geographic Information Society (CaGIS) Doctoral Scholarship Award, \$1,000.
  \item \dt{2026} IEEE GRSS Travel Grant --- IGARSS 2026, Washington, D.C., \$500.
  \item \dt{2026} NSF I-GUIDE Summer School: Spatial AI \& Convergence Science --- Selected Travel Award.
  \item \dt{2026} AAG-GISSG Student Honors Paper Competition --- 2nd Place, \$350.
  \item \dt{2026} AAG-IGIF Scholarship Award, AAG International Geographic Information Funds (IGIF), \$1,200.
  \item \dt{2025} Environment and Sustainability Graduate Fellow Award, Texas A\&M University, \$2,500.
  \item \dt{2025} Travel Grant, Summer School on Cyberinfrastructure and Disaster Resilience, TAMU (funded by NSF), \$200.
  \item \dt{2025} Texas A\&M Institute of Data Science (TAMIDS) Student Ambassador Scholarship, Domain Data Science Track, TAMU, \$2,000.
  \item \dt{2025} CaGIS International Travel Grant --- 2025 ICC, Vancouver, Canada, \$2,100.
  \item \dt{2025} International Cartographic Conference (ICC), Vancouver, Canada --- \award{Best Student Paper Award}.
  \item \dt{2025} AAG-GISSG Student Honors Paper Competition --- Honorable Mention, \$200 (Top 5).
  \item \dt{2025} AAG Applied Geography Specialty Group --- Student Travel Award, \$196.
  \item \dt{2024} Travel Grant (Visiting), Department of Geography, University of South Carolina, \$400.
  \item \dt{2024} Lifetime Membership --- Nu Theta Chapter, Gamma Theta Upsilon (International Geographic Honor Society).
  \item \dt{2024} Los Angeles Geospatial Summit --- ArcGIS StoryMaps Competition: Most Suitably Applied Analysis Methodology.
  \item \dt{2020} National First Prize --- Computer Design Competition for Chinese College Students (Big Data Practice).
  \item \dt{2020} Third Prize --- 10th MathorCup College Mathematical Modeling Challenge.
  \item \dt{2020} Second Prize --- 13th ``Certification Cup'' Mathematics China Mathematical Modeling Network Challenge (Inner Mongolia Region).
  \item \dt{2020} Third Prize --- China--US Youth Creators Competition (Haikou Region).
  \item \dt{2020} Second Prize --- 3rd China Youth Cup National University Student Mathematical Modeling Competition.
  \item \dt{2020} Third Prize --- 6th National Mobile Internet Innovation Competition (South China Region).
  \item \dt{2020} Third Prize --- Innovation Group of the 2020 China University Computer Competition (Artificial Intelligence Track).
  \item \dt{2020} Third Prize --- Finalist, Hainan Selection Competition of the 4th China Creative Wings Innovation Competition.
  \item \dt{2020} Honorable Mention --- 2020 Hainan Free Trade Port Entrepreneurship Competition.
  \item \dt{2020} Silver Award --- Hainan Creative Group, 6th China International ``Internet+'' Student Innovation and Entrepreneurship Competition.
  \item \dt{2020} Bronze Prize --- Hainan Region, Challenge Cup Student Entrepreneurship Plan Competition.
  \item \dt{2020} Second Prize --- 14th iCAN International Innovation and Entrepreneurship Competition (South China Region).
  \item \dt{2020--2021} First-Class Comprehensive Scholarship --- Hainan University, \textyen 2,500.
  \item \dt{2020--2021} Recognized as ``College Student with the Most Innovative Spirit and Practical Ability'' --- Hainan University.
\end{plainlist}

\vspace{0.4em}
\entry{2024--2025 NSF I-GUIDE Spatial AI Challenge --- Selected Abstract (2 consecutive years)}{}
\begin{tight}
  \item Selected in a competitive international challenge on Spatial AI and GeoAI innovation
  \item \textbf{2025:} GeoAgent4Disaster --- Autonomous multi-agent framework for multimodal disaster assessment
  \item \textbf{2024:} Perceiving Multidimensional Disaster Damages from Street-View Images Using Visual-Language Models
\end{tight}

%==============================================================
\section{Professional Activities and Service}
%==============================================================

\begin{plainlist}
  \item \dt{2025--2026} Data Science Student Ambassador, Texas A\&M Institute of Data Science (TAMIDS).
  \item \dt{2025--2026} Executive Committee Member, Aggie Geographers Graduate Society (AGGS).
  \item \dt{2025--2027} Student Co-Director of Remote Sensing Specialty Group, American Association of Geographers.
  \item \dt{2025--2027} Student Co-Director of Hazards, Risks, and Disasters, American Association of Geographers.
  \item \dt{2024--2025} Board Member, GISphere (responsible for the GISalon Roundtable Series).
  \item \dt{2021--2022} Vice President, Association of Robotics and Artificial Intelligence, Hainan University.
\end{plainlist}

\vspace{0.4em}
\textbf{Reviewer for leading journals and conferences, including:}
ACM Transactions on Autonomous and Adaptive Systems \,$\cdot$\, Computational Urban Science \,$\cdot$\, Climatic Change \,$\cdot$\, International Journal of Applied Earth Observation and Geoinformation \,$\cdot$\, International Journal of Geographical Information Science \,$\cdot$\, International Journal of ITS Research \,$\cdot$\, Scientific Reports \,$\cdot$\, Transactions in GIS \,$\cdot$\, Tourism Geographies \,$\cdot$\, Journal of Transport Geography \,$\cdot$\, International Journal of Remote Sensing \,$\cdot$\, 32nd SIGKDD Conference on Knowledge Discovery and Data Mining (KDD 2026) --- AI for Sciences Track

%==============================================================
\section{Open-Source Leadership}
%==============================================================

\entry{Founder, AutonomousGeoAI4Science (AutoGeoAI4Sci)}{\href{https://github.com/AutoGeoAI4Sci}{GitHub Organization}}
\vspace{0.15em}
Founded and led an open research community advancing autonomous and multimodal GeoAI for scientific discovery. The initiative integrates tutorials, open-source research code, cross-view modeling projects, and agent-based spatial intelligence systems.
\begin{tight}
  \item Curate and maintain 40+ repositories on cross-view GeoAI, disaster intelligence, and autonomous spatial reasoning
  \item Develop open-source tutorials and research frameworks for multimodal geospatial AI
  \item Build a collaborative platform connecting AI, GIScience, and Earth observation communities
  \item Promote reproducible research in autonomous geospatial intelligence
\end{tight}

%==============================================================
\section{Technical Skills}
%==============================================================

\begin{plainlist}
  \item \textbf{Applications:} PyTorch, TensorFlow, Scikit-learn, HuggingFace Transformers, Stable Diffusion, Visual Studio Code, Git/GitHub, Jupyter Notebook, Anaconda, Linux/Unix, Docker, Esri ArcGIS Pro, MS SQL Server, SQL Server Management Studio, PgAdmin, PostgreSQL \& PostGIS
  \item \textbf{Programming:} C, C++, Java, Python, SQL, R, ArcGIS API (JavaScript, Python), JavaScript, HTML
\end{plainlist}

\end{document}
